% ==========================================================================
\chapter{结论与展望}
% ==========================================================================

\section{主要结论}

本文以常微分方程初值问题的数值算法为主线，研究了粘性 Burgers 方程经有限体积离散后得到的刚性初值问题，系统分析了参数化 IMEX-Euler 分裂族的耦合刚性稳定性，并把相关语言推广到 Jin--Xin 松弛系统的源项分裂。主要结论可以概括为以下五点。

第一，建立了一个覆盖全显式到全隐式的单参数分裂族。通过引入 $\theta\in[0,1]$，把耗散算子按 $\theta$ 与 $1-\theta$ 的比例分配至显式与隐式部分。该族把全显式（$\theta=1$）、对半分裂（$\theta=1/2$）与全隐式（$\theta=0$）纳入同一框架，为系统比较不同分裂提供了设计空间。

第二，给出了耦合刚性稳定性的精确条件。对纯耗散情形，双变量稳定性函数沿刚性射线的值为
\begin{equation}
  R\left(-\theta x,-(1-\theta)x\right)=\frac{1-\theta x}{1+(1-\theta)x},
\end{equation}
稳定性 $|R|\le1$ 当且仅当
\begin{equation}
  (1-2\theta)x+2\ge0.
\end{equation}
因此，$\theta\le1/2$ 时格式对所有刚性稳定；$\theta>1/2$ 时仅对 $x\le2/(2\theta-1)$ 稳定。在该分裂族与纯耗散刚性极限下，这一条件既是必要的，也是充分的，且证明只用到初等代数与复平面圆盘的几何直观。对连续符号，进一步得到联合稳定条件 $\mu^2\le2d+(1-2\theta)d^2$；用刚性比 $\varepsilon=z_I/z_E$ 可把它写成 $|\varepsilon|\ge1$，而分裂族对应 $\varepsilon=(1-\theta)/\theta$。

第三，澄清了 A-稳定与 L-稳定的区别。沿刚性射线，只有 $\theta=0$ 是 L-稳定的（$R\to0$）；$0<\theta\le1/2$ 沿刚性射线稳定，但极限放大因子为 $-\theta/(1-\theta)\ne0$；$\theta=1/2$ 是临界情形，极限为 $-1$，表现为不衰减的锯齿振荡。因此，“稳定”不等于“刚性衰减”；对于需要快速压掉快模态的问题，应优先选择 $\theta=0$。

第四，完成了离散符号分析与数值验证。对一阶 Rusanov 格式，导出了网格尺度条件
\begin{equation}
  c+2(2\theta-1)r\le1\qquad(\theta>1/2),
\end{equation}
它解释了 $\theta>1/2$ 时最大步长按 $\Delta x^2/[2\nu(2\theta-1)]$ 收缩的现象。数值实验表明：刚性相图的边界与 $\theta=1/2$ 一致；$\theta=0.55$ 至 $1$ 的实测最大步长与线性预测相差不超过 $2\%$；MUSCL 格式的激波误差比一阶格式低 $15$ 至 $28$ 倍，观测收敛阶约为 $2$。

第五，把框架推广到 Jin--Xin 松弛系统。比较了源项全隐式（策略 A）与仅隐式化线性刚性部分（策略 B）两种分裂。谱分析表明，A 在 $\CFL\le1$ 时稳定，B 的临界 CFL 为 $1/3$，且该限制在 $\mu=0.1$ 到 $10^{4}$ 的范围内几乎不变。数值实验验证了 B 在 $\CFL=0.7$ 时的失稳，并揭示了 B 的平衡残差随 CFL 迅速增大的现象。

\section{局限性}

本文的分析存在以下局限，需要在后续研究中改进。

其一，稳定性分析基于线性化与冻结系数。对非线性 Burgers 方程，激波附近的系数变化、限制器的开关、边界条件等都可能影响稳定性。线性分析给出的是必要条件，非线性稳定性需要借助能量方法、TVD 性质或数值实验来补充。本文的数值实验虽然与理论一致，但没有给出严格的非线性稳定性证明。

其二，空间离散是一阶（稳定性分析）或二阶 MUSCL（精度实验），时间推进是一阶 IMEX-Euler。高阶 IMEX Runge--Kutta 方法、高阶 WENO 重构的稳定性边界可能不同；本文的结论不能直接推广到任意高阶格式。不过，“双变量稳定性函数 + 刚性射线”的框架是可迁移的，只需要把 $R(z_E,z_I)$ 换成相应格式的稳定性函数。

其三，研究对象是一维标量方程。对二维/三维系统、方程组、具有多个刚性参数的问题，稳定性函数将变成矩阵值，分析复杂度显著增加。本文的 Jin--Xin 部分已经展示了 $2\times2$ 矩阵的情形，但更一般的系统仍需进一步研究。

其四，$\theta$ 在本文中是常数。若让 $\theta$ 随时间或空间变化（自适应分裂），稳定性分析会变得更加困难；但这也可能带来更好的精度—效率平衡，是一个值得探索的方向。

其五，Jin--Xin 部分只考虑了无粘情形。粘性松弛系统的双刚性（扩散刚性 + 源项刚性）分析尚未展开；本文只给出了推广的思路，没有给出完整的定理与实验。

\section{展望}

基于本文的结果，可以在以下几个方向继续深入。

（1）\textbf{高阶格式的稳定性分析。} 把 IMEX-Euler 换成 IMEX Runge--Kutta 或 IMEX 线性多步格式，推导相应的双变量稳定性函数，研究阶数与刚性稳定性之间的权衡。特别是，能否构造在 $\theta>1/2$ 时仍保持 A-稳定的高阶格式？这需要新的分裂设计。

（2）\textbf{自适应分裂。} 让 $\theta$ 依赖于局部刚性 $|z_D|$：在刚性强的区域取较小的 $\theta$（更多隐式），在刚性弱的区域取较大的 $\theta$（更多显式）。这样可以在保证稳定性的同时降低隐式求解的规模。关键是如何设计 $\theta$ 的判据，使其既稳定又不会引入过多的数值耗散。

（3）\textbf{多维系统与算子分裂。} 把本文的框架推广到二维 Euler 方程、浅水方程等。此时扩散/粘性算子可能是张量型的，隐式求解需要多维线性系统；稳定性函数变成矩阵值，刚性射线的方向也更复杂。

（4）\textbf{非线性稳定性理论。} 用能量方法或相对熵方法证明 IMEX 格式在激波附近的有界性，而不仅仅依赖线性化。这将为“容许分裂”提供更坚实的理论基础。

（5）\textbf{与其他刚性处理方法的比较。} 除了 IMEX，还有指数积分、隐式—显式龙格—库塔、Rosenbrock 方法等。把本文的双变量稳定性语言与这些方法比较，可以更全面地理解“刚性分裂”的共性与差异。

（6）\textbf{软件实现与教学。} 本文的代码结构清晰、依赖少，适合作为数值 PDE 课程的教学案例。后续可以把它整理为一个开源教学包，包含更丰富的算例、可视化与自动稳定性分析工具。

总之，本文的工作是一个起点：它把“容许分裂”“双变量稳定性函数”“耦合刚性极限”“L-稳定”这套语言，落实为一个可以计算、可以作图、可以用数值实验检验的流程。希望这一流程能够为更复杂的双曲守恒律与刚性松弛系统的研究提供帮助。
